% !TEX program = xelatex
\def\StudyRoot{../..}
\def\StudyBook{MultipleChoice}
\documentclass[11pt,lang=cn,chinesefont=nofont,mode=simple,math=newtx,color=blue]{elegantbook}
\usepackage{longtable,array,fvextra,xurl,adjustbox,needspace}
\usepackage[normalem]{ulem}
\providecommand{\lt}{<}
\providecommand{\gt}{>}
\tracinglostchars=3
\setCJKmainfont{FandolSong-Regular.otf}[BoldFont=FandolSong-Bold.otf,ItalicFont=FandolKai-Regular.otf]
\setCJKsansfont{FandolHei-Regular.otf}[BoldFont=FandolHei-Bold.otf]
\setCJKmonofont{FandolFang-Regular.otf}
\setmonofont{DejaVuSansMono.ttf}[BoldFont=DejaVuSansMono-Bold.ttf,ItalicFont=DejaVuSansMono-Oblique.ttf,Scale=.84]
\geometry{a4paper,left=22mm,right=18mm,top=20mm,bottom=21mm,headheight=15pt,headsep=7mm}
\definecolor{structurecolor}{HTML}{173F57}
\colorlet{main}{structurecolor}
\colorlet{second}{structurecolor}
\colorlet{third}{structurecolor}
\hypersetup{unicode=true,colorlinks=true,linkcolor=structurecolor,urlcolor=structurecolor,citecolor=structurecolor,pdfauthor={CityU StudyNotes}}
\makeatletter
\renewcommand{\tableofcontents}{%
  \hypersetup{linktoc=all,linkcolor=black}%
  \chapter*{\contentsname}%
  \StudyTarget{\StudyBook__contents}\StudyTarget{\StudyBook___1}%
  \@starttoc{toc}%
  \hypersetup{linkcolor=structurecolor}\clearpage}
\makeatother
\setstretch{1.25}
\setlength{\parskip}{.35\baselineskip}
\setlength{\parindent}{0pt}
\setlength{\emergencystretch}{2em}
\setcounter{tocdepth}{1}
\setcounter{secnumdepth}{-1}
\setlist{nosep,leftmargin=1.8em,topsep=.35\baselineskip}
\titleformat{\chapter}[display]{\Huge\sffamily\bfseries\color{structurecolor}}{}{0pt}{}[\vspace{.4ex}\titlerule]
\titlespacing*{\chapter}{0pt}{0pt}{22pt}
\titleformat{\section}{\Large\sffamily\bfseries\color{structurecolor}}{}{0pt}{}
\titleformat{\subsection}{\large\sffamily\bfseries}{}{0pt}{}
\titleformat{\subsubsection}{\normalsize\sffamily\bfseries}{}{0pt}{}
\titlespacing*{\section}{0pt}{1.6\baselineskip}{.5\baselineskip}
\titlespacing*{\subsection}{0pt}{1.1\baselineskip}{.35\baselineskip}
\providecommand{\StudyBook}{Sample}
\providecommand{\StudyRoot}{../..}
\providecommand{\StudySourceRef}{main}
\providecommand{\StudyDate}{2026-10-16}
\newcommand{\StudyRunning}{\StudyBook}
\fancyhf{}
\fancyhead[L]{\small\sffamily\color{gray}\StudyRunning}
\fancyhead[R]{\small\sffamily\color{gray}CityU StudyNotes}
\fancyfoot[R]{\small\thepage}
\renewcommand{\headrulewidth}{.3pt}
\pagestyle{fancy}
\fancypagestyle{plain}{\fancyhf{}\fancyfoot[R]{\small\thepage}\renewcommand{\headrulewidth}{0pt}}
\newwrite\StudyMap
\AtBeginDocument{\immediate\openout\StudyMap=\jobname.pages}
\makeatletter
\newcommand{\StudyTarget}[1]{%
  \hypertarget{#1}{}\label{#1}%
  \protected@write\StudyMap{}{#1|\thepage}%
}
\makeatother
\newcommand{\StudyChapter}[2]{%
  \clearpage\renewcommand{\StudyRunning}{\ifcsname studytitle@#1\endcsname\csname studytitle@#1\endcsname\else#1\fi}%
  \chapter*{#2}\phantomsection\addcontentsline{toc}{chapter}{#2}%
  \StudyTarget{#1__root}%
}
\newcommand{\StudyHeading}[2]{%
  \ifcase#1\relax\or
    \section*{#2}\phantomsection\addcontentsline{toc}{section}{#2}%
  \or\subsection*{#2}\or\subsubsection*{#2}\else\paragraph*{#2}\fi
}
\newcommand{\StudyImage}[1]{%
  \par\medskip\begingroup
  \edef\StudyResolvedImage{\ifcsname studyimage@#1\endcsname\csname studyimage@#1\endcsname\else\StudyRoot/#1\fi}%
  \noindent\makebox[\linewidth][c]{%
  \adjustbox{max width=.96\linewidth,max height=.69\textheight}{%
  \includegraphics{\StudyResolvedImage}}}\endgroup\par\medskip
}
\newcommand{\StudyCellImage}[1]{%
  \begingroup
  \edef\StudyResolvedImage{\ifcsname studyimage@#1\endcsname\csname studyimage@#1\endcsname\else\StudyRoot/#1\fi}%
  \adjustbox{max width=\linewidth,max height=.28\textheight}{\includegraphics{\StudyResolvedImage}}\endgroup}
\DefineVerbatimEnvironment{StudyCode}{Verbatim}{breaklines=true,breakanywhere=true,fontsize=\small,tabsize=4,framesep=3mm,frame=leftline,rulecolor=\color{gray!55}}
\newcommand{\StudyInline}[1]{\allowbreak{}\texttt{#1}\allowbreak{}}
\newenvironment{StudySolution}[1]{%
  \par\medskip\Needspace{3\baselineskip}\noindent
  {\bfseries\color{structurecolor}#1}\par\nobreak\smallskip
}{\par\medskip}
\newcommand{\StudyBadge}[2]{%
  \begingroup\small\sffamily
  \ifstrequal{#1}{mid}{\colorlet{studybadge}{green!35!black}}{%
  \ifstrequal{#1}{final}{\colorlet{studybadge}{yellow!45!black}}{\colorlet{studybadge}{red!45!black}}}%
  \textcolor{studybadge}{\rule[-.2ex]{1.2pt}{1.2em}}\enspace\textcolor{black}{#2}\quad\endgroup}
\newcommand{\StudyRef}[3]{%
  \ifcsname studyref@#1@#2\endcsname
    \csname studyref@#1@#2\endcsname{#3}%
  \else #3\textsuperscript{[待定页]}\fi}
\newcommand{\StudyRepoRef}[3]{%
  \href{https://github.com/CrazyShout/CityU-StudyNotes/blob/\StudySourceRef/#1\ifstrempty{#2}{}{\##2}}{#3}}
\newcommand{\StudyNotebookCell}[2]{%
  \ifcsname studynb@#1@#2\endcsname
    \input{\csname studynb@#1@#2\endcsname}%
  \else\PackageError{studynotes}{Missing saved notebook cell #2}{}\fi}
\newcommand{\StudyCoverExtra}{}
\newcommand{\StudyCover}[2]{%
  \thispagestyle{empty}
  \StudyTarget{\StudyBook__root}\StudyCoverExtra
  \vspace*{22mm}
  {\sffamily\large\color{structurecolor}CITYU STUDYNOTES\par}
  \vspace{12mm}
  {\sffamily\bfseries\Huge\color{structurecolor}#1\par}
  \vspace{6mm}{\Large #2\par}
  \vspace{10mm}\noindent\textcolor{structurecolor}{\rule{.42\linewidth}{1pt}}\par
  \vspace{8mm}中文讲解 · English terminology · 双语题答\par
  A4 单栏 · 根据已下载教师材料整理\par
  \vfill
  {\small 本资料用于课程学习与复习，保留原材料、历史题源与教学补充的身份。\par
  模板：ElegantBook 4.7；版式与构建由 CityU-StudyNotes 维护。\par
  \StudyDate\par}
  \clearpage
}
\newcommand{\StudyQuestionExtra}{}
\newcommand{\StudyQuestion}[3]{%
  \Needspace{10\baselineskip}
  \section*{#1 · #2}\phantomsection\addcontentsline{toc}{section}{#1 · #2}%
  \StudyTarget{\StudyBook__#1}%
  \StudyQuestionExtra
  \lowercase{\def\StudyQuestionAnchor{#1}}%
  \noindent\begin{minipage}[t]{.65\linewidth}{\small\sffamily #3}\end{minipage}\hfill
  \begin{minipage}[t]{.33\linewidth}\raggedleft\scriptsize
  \StudyRef{Questions}{\StudyQuestionAnchor}{题目} · \StudyRef{Answers}{\StudyQuestionAnchor}{答案}
  \end{minipage}\par\medskip
}
\InputIfFileExists{\StudyRoot/.build/latex/\StudyBook/refs.tex}{}{}

\begin{document}
\setlength{\parskip}{.25\baselineskip}
\titlespacing*{\subsection}{0pt}{.7\baselineskip}{.2\baselineskip}
\renewcommand{\StudyCoverExtra}{\StudyTarget{MultipleChoice__cs5489-multiple-select-questions-and-answers}}
\StudyCover{CS5489 · 期中选择题与答案}{Multiple-Select Questions and Answers}
\input{\StudyRoot/CS5489/exam-review/Midterm/frontmatter/MultipleChoice.tex}
\begingroup\small\setlength{\cftbeforechapskip}{.45\baselineskip}\tableofcontents\endgroup
\Needspace{10\baselineskip}\phantomsection\addcontentsline{toc}{chapter}{学习流程与模型诊断}
\renewcommand{\StudyQuestionExtra}{\StudyTarget{MultipleChoice__mt004}\StudyTarget{MultipleChoice__mt004-cv}}
\StudyQuestion{MT004}{CV选完参数后，为什么还要重新训练？}{选择题（多选） / Multiple select}
\input{\StudyRoot/CS5489/exam-review/Midterm/items/MT004/provenance.tex}
\subsection*{English question}
\StudyTarget{MultipleChoice__english-question}
\input{\StudyRoot/CS5489/exam-review/Midterm/items/MT004/question-en.tex}
\subsection*{中文题意}
\StudyTarget{MultipleChoice___2}
\input{\StudyRoot/CS5489/exam-review/Midterm/items/MT004/question-zh.tex}
\subsection*{English answer}
\StudyTarget{MultipleChoice__english-answer}
\input{\StudyRoot/CS5489/exam-review/Midterm/items/MT004/answer-en.tex}
\subsection*{中文对应答案}
\StudyTarget{MultipleChoice___3}
\input{\StudyRoot/CS5489/exam-review/Midterm/items/MT004/answer-zh.tex}
\input{\StudyRoot/CS5489/exam-review/Midterm/items/MT004/concise-extra.tex}
\renewcommand{\StudyQuestionExtra}{\StudyTarget{MultipleChoice__mt028}}
\StudyQuestion{MT028}{训练误差低、测试误差高：可能原因}{选择题（多选） / Multiple select}
\input{\StudyRoot/CS5489/exam-review/Midterm/items/MT028/provenance.tex}
\subsection*{English question}
\StudyTarget{MultipleChoice__english-question_1}
\input{\StudyRoot/CS5489/exam-review/Midterm/items/MT028/question-en.tex}
\subsection*{中文题意}
\StudyTarget{MultipleChoice___4}
\input{\StudyRoot/CS5489/exam-review/Midterm/items/MT028/question-zh.tex}
\subsection*{English answer}
\StudyTarget{MultipleChoice__english-answer_1}
\input{\StudyRoot/CS5489/exam-review/Midterm/items/MT028/answer-en.tex}
\subsection*{中文对应答案}
\StudyTarget{MultipleChoice___5}
\input{\StudyRoot/CS5489/exam-review/Midterm/items/MT028/answer-zh.tex}
\input{\StudyRoot/CS5489/exam-review/Midterm/items/MT028/concise-extra.tex}
\renewcommand{\StudyQuestionExtra}{\StudyTarget{MultipleChoice__mt002}}
\StudyQuestion{MT002}{训练准、测试差：哪些改进值得尝试？}{选择题（多选） / Multiple select}
\input{\StudyRoot/CS5489/exam-review/Midterm/items/MT002/provenance.tex}
\subsection*{English question}
\StudyTarget{MultipleChoice__english-question_2}
\input{\StudyRoot/CS5489/exam-review/Midterm/items/MT002/question-en.tex}
\subsection*{中文题意}
\StudyTarget{MultipleChoice___6}
\input{\StudyRoot/CS5489/exam-review/Midterm/items/MT002/question-zh.tex}
\subsection*{English answer}
\StudyTarget{MultipleChoice__english-answer_2}
\input{\StudyRoot/CS5489/exam-review/Midterm/items/MT002/answer-en.tex}
\subsection*{中文对应答案}
\StudyTarget{MultipleChoice___7}
\input{\StudyRoot/CS5489/exam-review/Midterm/items/MT002/answer-zh.tex}
\input{\StudyRoot/CS5489/exam-review/Midterm/items/MT002/concise-extra.tex}
\Needspace{10\baselineskip}\phantomsection\addcontentsline{toc}{chapter}{Bayes、Naive Bayes与模型比较}
\renewcommand{\StudyQuestionExtra}{\StudyTarget{MultipleChoice__mt040}\StudyTarget{MultipleChoice__l2q05}\StudyTarget{MultipleChoice__mt040-bayesnaive-bayes}}
\StudyQuestion{MT040}{Bayes规则和Naive Bayes的关系}{选择题（多选） / Multiple select}
\input{\StudyRoot/CS5489/exam-review/Midterm/items/MT040/provenance.tex}
\subsection*{English question}
\StudyTarget{MultipleChoice__english-question_3}
\input{\StudyRoot/CS5489/exam-review/Midterm/items/MT040/question-en.tex}
\subsection*{中文题意}
\StudyTarget{MultipleChoice___8}
\input{\StudyRoot/CS5489/exam-review/Midterm/items/MT040/question-zh.tex}
\subsection*{English answer}
\StudyTarget{MultipleChoice__english-answer_3}
\input{\StudyRoot/CS5489/exam-review/Midterm/items/MT040/answer-en.tex}
\subsection*{中文对应答案}
\StudyTarget{MultipleChoice___9}
\input{\StudyRoot/CS5489/exam-review/Midterm/items/MT040/answer-zh.tex}
\input{\StudyRoot/CS5489/exam-review/Midterm/items/MT040/concise-extra.tex}
\renewcommand{\StudyQuestionExtra}{\StudyTarget{MultipleChoice__mt052}\StudyTarget{MultipleChoice__l2q01}}
\StudyQuestion{MT052}{生成式分类器学了什么}{选择题（多选） / Multiple select}
\input{\StudyRoot/CS5489/exam-review/Midterm/items/MT052/provenance.tex}
\subsection*{English question}
\StudyTarget{MultipleChoice__english-question_4}
\input{\StudyRoot/CS5489/exam-review/Midterm/items/MT052/question-en.tex}
\subsection*{中文题意}
\StudyTarget{MultipleChoice___10}
\input{\StudyRoot/CS5489/exam-review/Midterm/items/MT052/question-zh.tex}
\subsection*{English answer}
\StudyTarget{MultipleChoice__english-answer_4}
\input{\StudyRoot/CS5489/exam-review/Midterm/items/MT052/answer-en.tex}
\subsection*{中文对应答案}
\StudyTarget{MultipleChoice___11}
\input{\StudyRoot/CS5489/exam-review/Midterm/items/MT052/answer-zh.tex}
\input{\StudyRoot/CS5489/exam-review/Midterm/items/MT052/concise-extra.tex}
\renewcommand{\StudyQuestionExtra}{\StudyTarget{MultipleChoice__mt015}\StudyTarget{MultipleChoice__l2q08}\StudyTarget{MultipleChoice__mt015-nb}}
\StudyQuestion{MT015}{NB的最优性、独立性与边界}{选择题（多选） / Multiple select}
\input{\StudyRoot/CS5489/exam-review/Midterm/items/MT015/provenance.tex}
\subsection*{English question}
\StudyTarget{MultipleChoice__english-question_5}
\input{\StudyRoot/CS5489/exam-review/Midterm/items/MT015/question-en.tex}
\subsection*{中文题意}
\StudyTarget{MultipleChoice___12}
\input{\StudyRoot/CS5489/exam-review/Midterm/items/MT015/question-zh.tex}
\subsection*{English answer}
\StudyTarget{MultipleChoice__english-answer_5}
\input{\StudyRoot/CS5489/exam-review/Midterm/items/MT015/answer-en.tex}
\subsection*{中文对应答案}
\StudyTarget{MultipleChoice___13}
\input{\StudyRoot/CS5489/exam-review/Midterm/items/MT015/answer-zh.tex}
\input{\StudyRoot/CS5489/exam-review/Midterm/items/MT015/concise-extra.tex}
\renewcommand{\StudyQuestionExtra}{\StudyTarget{MultipleChoice__mt032}\StudyTarget{MultipleChoice__l2q06}\StudyTarget{MultipleChoice__mt032-bayesian-classifier}}
\StudyQuestion{MT032}{Bayesian classifier的性质}{选择题（多选） / Multiple select}
\input{\StudyRoot/CS5489/exam-review/Midterm/items/MT032/provenance.tex}
\subsection*{English question}
\StudyTarget{MultipleChoice__english-question_6}
\input{\StudyRoot/CS5489/exam-review/Midterm/items/MT032/question-en.tex}
\subsection*{中文题意}
\StudyTarget{MultipleChoice___14}
\input{\StudyRoot/CS5489/exam-review/Midterm/items/MT032/question-zh.tex}
\subsection*{English answer}
\StudyTarget{MultipleChoice__english-answer_6}
\input{\StudyRoot/CS5489/exam-review/Midterm/items/MT032/answer-en.tex}
\subsection*{中文对应答案}
\StudyTarget{MultipleChoice___15}
\input{\StudyRoot/CS5489/exam-review/Midterm/items/MT032/answer-zh.tex}
\input{\StudyRoot/CS5489/exam-review/Midterm/items/MT032/concise-extra.tex}
\renewcommand{\StudyQuestionExtra}{\StudyTarget{MultipleChoice__mt005}\StudyTarget{MultipleChoice__l2q07}\StudyTarget{MultipleChoice__mt005-nb}}
\StudyQuestion{MT005}{NB的学习、不确定性和正则化}{选择题（多选） / Multiple select}
\input{\StudyRoot/CS5489/exam-review/Midterm/items/MT005/provenance.tex}
\subsection*{English question}
\StudyTarget{MultipleChoice__english-question_7}
\input{\StudyRoot/CS5489/exam-review/Midterm/items/MT005/question-en.tex}
\subsection*{中文题意}
\StudyTarget{MultipleChoice___16}
\input{\StudyRoot/CS5489/exam-review/Midterm/items/MT005/question-zh.tex}
\subsection*{English answer}
\StudyTarget{MultipleChoice__english-answer_7}
\input{\StudyRoot/CS5489/exam-review/Midterm/items/MT005/answer-en.tex}
\subsection*{中文对应答案}
\StudyTarget{MultipleChoice___17}
\input{\StudyRoot/CS5489/exam-review/Midterm/items/MT005/answer-zh.tex}
\input{\StudyRoot/CS5489/exam-review/Midterm/items/MT005/concise-extra.tex}
\renewcommand{\StudyQuestionExtra}{\StudyTarget{MultipleChoice__mt014}\StudyTarget{MultipleChoice__l2q03}}
\StudyQuestion{MT014}{生成式与判别式的区别}{选择题（多选） / Multiple select · 选不正确项 / Select incorrect statements}
\input{\StudyRoot/CS5489/exam-review/Midterm/items/MT014/provenance.tex}
\subsection*{English question}
\StudyTarget{MultipleChoice__english-question_8}
\input{\StudyRoot/CS5489/exam-review/Midterm/items/MT014/question-en.tex}
\subsection*{中文题意}
\StudyTarget{MultipleChoice___18}
\input{\StudyRoot/CS5489/exam-review/Midterm/items/MT014/question-zh.tex}
\subsection*{English answer}
\StudyTarget{MultipleChoice__english-answer_8}
\input{\StudyRoot/CS5489/exam-review/Midterm/items/MT014/answer-en.tex}
\subsection*{中文对应答案}
\StudyTarget{MultipleChoice___19}
\input{\StudyRoot/CS5489/exam-review/Midterm/items/MT014/answer-zh.tex}
\input{\StudyRoot/CS5489/exam-review/Midterm/items/MT014/concise-extra.tex}
\renewcommand{\StudyQuestionExtra}{\StudyTarget{MultipleChoice__mt017}\StudyTarget{MultipleChoice__l2q11}}
\StudyQuestion{MT017}{哪些模型能产生不同形状的边界}{选择题（多选） / Multiple select}
\input{\StudyRoot/CS5489/exam-review/Midterm/items/MT017/provenance.tex}
\subsection*{English question}
\StudyTarget{MultipleChoice__english-question_9}
\input{\StudyRoot/CS5489/exam-review/Midterm/items/MT017/question-en.tex}
\subsection*{中文题意}
\StudyTarget{MultipleChoice___20}
\input{\StudyRoot/CS5489/exam-review/Midterm/items/MT017/question-zh.tex}
\subsection*{English answer}
\StudyTarget{MultipleChoice__english-answer_9}
\input{\StudyRoot/CS5489/exam-review/Midterm/items/MT017/answer-en.tex}
\subsection*{中文对应答案}
\StudyTarget{MultipleChoice___21}
\input{\StudyRoot/CS5489/exam-review/Midterm/items/MT017/answer-zh.tex}
\input{\StudyRoot/CS5489/exam-review/Midterm/items/MT017/concise-extra.tex}
\renewcommand{\StudyQuestionExtra}{\StudyTarget{MultipleChoice__mt027}\StudyTarget{MultipleChoice__l2q10}\StudyTarget{MultipleChoice__mt027-nb}}
\StudyQuestion{MT027}{连续特征能否使用NB：混合选择题}{选择题（多选） / Multiple select}
\input{\StudyRoot/CS5489/exam-review/Midterm/items/MT027/provenance.tex}
\subsection*{English question}
\StudyTarget{MultipleChoice__english-question_10}
\input{\StudyRoot/CS5489/exam-review/Midterm/items/MT027/question-en.tex}
\subsection*{中文题意}
\StudyTarget{MultipleChoice___22}
\input{\StudyRoot/CS5489/exam-review/Midterm/items/MT027/question-zh.tex}
\subsection*{English answer}
\StudyTarget{MultipleChoice__english-answer_10}
\input{\StudyRoot/CS5489/exam-review/Midterm/items/MT027/answer-en.tex}
\subsection*{中文对应答案}
\StudyTarget{MultipleChoice___23}
\input{\StudyRoot/CS5489/exam-review/Midterm/items/MT027/answer-zh.tex}
\input{\StudyRoot/CS5489/exam-review/Midterm/items/MT027/concise-extra.tex}
\Needspace{10\baselineskip}\phantomsection\addcontentsline{toc}{chapter}{逻辑回归与参数正则}
\renewcommand{\StudyQuestionExtra}{\StudyTarget{MultipleChoice__mt041}\StudyTarget{MultipleChoice__mt041-logistic}}
\StudyQuestion{MT041}{Logistic的概率、多分类与最优解}{选择题（多选） / Multiple select}
\input{\StudyRoot/CS5489/exam-review/Midterm/items/MT041/provenance.tex}
\subsection*{English question}
\StudyTarget{MultipleChoice__english-question_11}
\input{\StudyRoot/CS5489/exam-review/Midterm/items/MT041/question-en.tex}
\subsection*{中文题意}
\StudyTarget{MultipleChoice___24}
\input{\StudyRoot/CS5489/exam-review/Midterm/items/MT041/question-zh.tex}
\subsection*{English answer}
\StudyTarget{MultipleChoice__english-answer_11}
\input{\StudyRoot/CS5489/exam-review/Midterm/items/MT041/answer-en.tex}
\subsection*{中文对应答案}
\StudyTarget{MultipleChoice___25}
\input{\StudyRoot/CS5489/exam-review/Midterm/items/MT041/answer-zh.tex}
\input{\StudyRoot/CS5489/exam-review/Midterm/items/MT041/concise-extra.tex}
\renewcommand{\StudyQuestionExtra}{\StudyTarget{MultipleChoice__mt018}\StudyTarget{MultipleChoice__mt018-l2}}
\StudyQuestion{MT018}{为什么在线性分类器上加L2？}{选择题（多选） / Multiple select}
\input{\StudyRoot/CS5489/exam-review/Midterm/items/MT018/provenance.tex}
\subsection*{English question}
\StudyTarget{MultipleChoice__english-question_12}
\input{\StudyRoot/CS5489/exam-review/Midterm/items/MT018/question-en.tex}
\subsection*{中文题意}
\StudyTarget{MultipleChoice___26}
\input{\StudyRoot/CS5489/exam-review/Midterm/items/MT018/question-zh.tex}
\subsection*{English answer}
\StudyTarget{MultipleChoice__english-answer_12}
\input{\StudyRoot/CS5489/exam-review/Midterm/items/MT018/answer-en.tex}
\subsection*{中文对应答案}
\StudyTarget{MultipleChoice___27}
\input{\StudyRoot/CS5489/exam-review/Midterm/items/MT018/answer-zh.tex}
\input{\StudyRoot/CS5489/exam-review/Midterm/items/MT018/concise-extra.tex}
\renewcommand{\StudyQuestionExtra}{\StudyTarget{MultipleChoice__mt001}\StudyTarget{MultipleChoice__mt001-logistic}}
\StudyQuestion{MT001}{Logistic训练：哪些说法不成立？}{选择题（多选） / Multiple select · 选不正确项 / Select incorrect statements}
\input{\StudyRoot/CS5489/exam-review/Midterm/items/MT001/provenance.tex}
\subsection*{English question}
\StudyTarget{MultipleChoice__english-question_13}
\input{\StudyRoot/CS5489/exam-review/Midterm/items/MT001/question-en.tex}
\subsection*{中文题意}
\StudyTarget{MultipleChoice___28}
\input{\StudyRoot/CS5489/exam-review/Midterm/items/MT001/question-zh.tex}
\subsection*{English answer}
\StudyTarget{MultipleChoice__english-answer_13}
\input{\StudyRoot/CS5489/exam-review/Midterm/items/MT001/answer-en.tex}
\subsection*{中文对应答案}
\StudyTarget{MultipleChoice___29}
\input{\StudyRoot/CS5489/exam-review/Midterm/items/MT001/answer-zh.tex}
\input{\StudyRoot/CS5489/exam-review/Midterm/items/MT001/concise-extra.tex}
\renewcommand{\StudyQuestionExtra}{\StudyTarget{MultipleChoice__mt053}\StudyTarget{MultipleChoice__l2q14}\StudyTarget{MultipleChoice__mt053-lr}}
\StudyQuestion{MT053}{LR与生成式模型：混合选择题}{选择题（多选） / Multiple select}
\input{\StudyRoot/CS5489/exam-review/Midterm/items/MT053/provenance.tex}
\subsection*{English question}
\StudyTarget{MultipleChoice__english-question_14}
\input{\StudyRoot/CS5489/exam-review/Midterm/items/MT053/question-en.tex}
\subsection*{中文题意}
\StudyTarget{MultipleChoice___30}
\input{\StudyRoot/CS5489/exam-review/Midterm/items/MT053/question-zh.tex}
\subsection*{English answer}
\StudyTarget{MultipleChoice__english-answer_14}
\input{\StudyRoot/CS5489/exam-review/Midterm/items/MT053/answer-en.tex}
\subsection*{中文对应答案}
\StudyTarget{MultipleChoice___31}
\input{\StudyRoot/CS5489/exam-review/Midterm/items/MT053/answer-zh.tex}
\input{\StudyRoot/CS5489/exam-review/Midterm/items/MT053/concise-extra.tex}
\Needspace{10\baselineskip}\phantomsection\addcontentsline{toc}{chapter}{SVM与核方法}
\renewcommand{\StudyQuestionExtra}{\StudyTarget{MultipleChoice__mt042}\StudyTarget{MultipleChoice__mt042-svm}}
\StudyQuestion{MT042}{SVM的基本特点}{选择题（多选） / Multiple select}
\input{\StudyRoot/CS5489/exam-review/Midterm/items/MT042/provenance.tex}
\subsection*{English question}
\StudyTarget{MultipleChoice__english-question_15}
\input{\StudyRoot/CS5489/exam-review/Midterm/items/MT042/question-en.tex}
\subsection*{中文题意}
\StudyTarget{MultipleChoice___32}
\input{\StudyRoot/CS5489/exam-review/Midterm/items/MT042/question-zh.tex}
\subsection*{English answer}
\StudyTarget{MultipleChoice__english-answer_15}
\input{\StudyRoot/CS5489/exam-review/Midterm/items/MT042/answer-en.tex}
\subsection*{中文对应答案}
\StudyTarget{MultipleChoice___33}
\input{\StudyRoot/CS5489/exam-review/Midterm/items/MT042/answer-zh.tex}
\input{\StudyRoot/CS5489/exam-review/Midterm/items/MT042/concise-extra.tex}
\renewcommand{\StudyQuestionExtra}{\StudyTarget{MultipleChoice__mt016}\StudyTarget{MultipleChoice__mt016-svm}}
\StudyQuestion{MT016}{核SVM：变换、存储与非向量输入}{选择题（多选） / Multiple select}
\input{\StudyRoot/CS5489/exam-review/Midterm/items/MT016/provenance.tex}
\subsection*{English question}
\StudyTarget{MultipleChoice__english-question_16}
\input{\StudyRoot/CS5489/exam-review/Midterm/items/MT016/question-en.tex}
\subsection*{中文题意}
\StudyTarget{MultipleChoice___34}
\input{\StudyRoot/CS5489/exam-review/Midterm/items/MT016/question-zh.tex}
\subsection*{English answer}
\StudyTarget{MultipleChoice__english-answer_16}
\input{\StudyRoot/CS5489/exam-review/Midterm/items/MT016/answer-en.tex}
\subsection*{中文对应答案}
\StudyTarget{MultipleChoice___35}
\input{\StudyRoot/CS5489/exam-review/Midterm/items/MT016/answer-zh.tex}
\input{\StudyRoot/CS5489/exam-review/Midterm/items/MT016/concise-extra.tex}
\renewcommand{\StudyQuestionExtra}{\StudyTarget{MultipleChoice__mt031}\StudyTarget{MultipleChoice__mt031-rbf}}
\StudyQuestion{MT031}{RBF距离与带宽：看近点和远点}{选择题（多选） / Multiple select}
\input{\StudyRoot/CS5489/exam-review/Midterm/items/MT031/provenance.tex}
\subsection*{English question}
\StudyTarget{MultipleChoice__english-question_17}
\input{\StudyRoot/CS5489/exam-review/Midterm/items/MT031/question-en.tex}
\subsection*{中文题意}
\StudyTarget{MultipleChoice___36}
\input{\StudyRoot/CS5489/exam-review/Midterm/items/MT031/question-zh.tex}
\subsection*{English answer}
\StudyTarget{MultipleChoice__english-answer_17}
\input{\StudyRoot/CS5489/exam-review/Midterm/items/MT031/answer-en.tex}
\subsection*{中文对应答案}
\StudyTarget{MultipleChoice___37}
\input{\StudyRoot/CS5489/exam-review/Midterm/items/MT031/answer-zh.tex}
\input{\StudyRoot/CS5489/exam-review/Midterm/items/MT031/concise-extra.tex}
\renewcommand{\StudyQuestionExtra}{\StudyTarget{MultipleChoice__mt055}}
\StudyQuestion{MT055}{判断五个候选核是否合法}{选择题（多选） / Multiple select}
\input{\StudyRoot/CS5489/exam-review/Midterm/items/MT055/provenance.tex}
\subsection*{English question}
\StudyTarget{MultipleChoice__english-question_18}
\input{\StudyRoot/CS5489/exam-review/Midterm/items/MT055/question-en.tex}
\subsection*{中文题意}
\StudyTarget{MultipleChoice___38}
\input{\StudyRoot/CS5489/exam-review/Midterm/items/MT055/question-zh.tex}
\subsection*{English answer}
\StudyTarget{MultipleChoice__english-answer_18}
\input{\StudyRoot/CS5489/exam-review/Midterm/items/MT055/answer-en.tex}
\subsection*{中文对应答案}
\StudyTarget{MultipleChoice___39}
\input{\StudyRoot/CS5489/exam-review/Midterm/items/MT055/answer-zh.tex}
\input{\StudyRoot/CS5489/exam-review/Midterm/items/MT055/concise-extra.tex}
\renewcommand{\StudyQuestionExtra}{\StudyTarget{MultipleChoice__mt054}\StudyTarget{MultipleChoice__mt054-crbf}}
\StudyQuestion{MT054}{C、支持向量与RBF的零训练误差}{选择题（多选） / Multiple select}
\input{\StudyRoot/CS5489/exam-review/Midterm/items/MT054/provenance.tex}
\subsection*{English question}
\StudyTarget{MultipleChoice__english-question_19}
\input{\StudyRoot/CS5489/exam-review/Midterm/items/MT054/question-en.tex}
\subsection*{中文题意}
\StudyTarget{MultipleChoice___40}
\input{\StudyRoot/CS5489/exam-review/Midterm/items/MT054/question-zh.tex}
\subsection*{English answer}
\StudyTarget{MultipleChoice__english-answer_19}
\input{\StudyRoot/CS5489/exam-review/Midterm/items/MT054/answer-en.tex}
\subsection*{中文对应答案}
\StudyTarget{MultipleChoice___41}
\input{\StudyRoot/CS5489/exam-review/Midterm/items/MT054/answer-zh.tex}
\input{\StudyRoot/CS5489/exam-review/Midterm/items/MT054/concise-extra.tex}
\renewcommand{\StudyQuestionExtra}{\StudyTarget{MultipleChoice__mt043}\StudyTarget{MultipleChoice__mt043-svm}}
\StudyQuestion{MT043}{核SVM怎样处理非线性？}{选择题（多选） / Multiple select}
\input{\StudyRoot/CS5489/exam-review/Midterm/items/MT043/provenance.tex}
\subsection*{English question}
\StudyTarget{MultipleChoice__english-question_20}
\input{\StudyRoot/CS5489/exam-review/Midterm/items/MT043/question-en.tex}
\subsection*{中文题意}
\StudyTarget{MultipleChoice___42}
\input{\StudyRoot/CS5489/exam-review/Midterm/items/MT043/question-zh.tex}
\subsection*{English answer}
\StudyTarget{MultipleChoice__english-answer_20}
\input{\StudyRoot/CS5489/exam-review/Midterm/items/MT043/answer-en.tex}
\subsection*{中文对应答案}
\StudyTarget{MultipleChoice___43}
\input{\StudyRoot/CS5489/exam-review/Midterm/items/MT043/answer-zh.tex}
\input{\StudyRoot/CS5489/exam-review/Midterm/items/MT043/concise-extra.tex}
\renewcommand{\StudyQuestionExtra}{\StudyTarget{MultipleChoice__mt029}\StudyTarget{MultipleChoice__mt029-svm}}
\StudyQuestion{MT029}{哪些SVM说法错误？}{选择题（多选） / Multiple select · 选不正确项 / Select incorrect statements}
\input{\StudyRoot/CS5489/exam-review/Midterm/items/MT029/provenance.tex}
\subsection*{English question}
\StudyTarget{MultipleChoice__english-question_21}
\input{\StudyRoot/CS5489/exam-review/Midterm/items/MT029/question-en.tex}
\subsection*{中文题意}
\StudyTarget{MultipleChoice___44}
\input{\StudyRoot/CS5489/exam-review/Midterm/items/MT029/question-zh.tex}
\subsection*{English answer}
\StudyTarget{MultipleChoice__english-answer_21}
\input{\StudyRoot/CS5489/exam-review/Midterm/items/MT029/answer-en.tex}
\subsection*{中文对应答案}
\StudyTarget{MultipleChoice___45}
\input{\StudyRoot/CS5489/exam-review/Midterm/items/MT029/answer-zh.tex}
\input{\StudyRoot/CS5489/exam-review/Midterm/items/MT029/concise-extra.tex}
\Needspace{10\baselineskip}\phantomsection\addcontentsline{toc}{chapter}{回归、特征选择与残差损失}
\renewcommand{\StudyQuestionExtra}{\StudyTarget{MultipleChoice__mt019}\StudyTarget{MultipleChoice__mt019-ridgelasso}}
\StudyQuestion{MT019}{Ridge、LASSO与零系数}{选择题（多选） / Multiple select}
\input{\StudyRoot/CS5489/exam-review/Midterm/items/MT019/provenance.tex}
\subsection*{English question}
\StudyTarget{MultipleChoice__english-question_22}
\input{\StudyRoot/CS5489/exam-review/Midterm/items/MT019/question-en.tex}
\subsection*{中文题意}
\StudyTarget{MultipleChoice___46}
\input{\StudyRoot/CS5489/exam-review/Midterm/items/MT019/question-zh.tex}
\subsection*{English answer}
\StudyTarget{MultipleChoice__english-answer_22}
\input{\StudyRoot/CS5489/exam-review/Midterm/items/MT019/answer-en.tex}
\subsection*{中文对应答案}
\StudyTarget{MultipleChoice___47}
\input{\StudyRoot/CS5489/exam-review/Midterm/items/MT019/answer-zh.tex}
\input{\StudyRoot/CS5489/exam-review/Midterm/items/MT019/concise-extra.tex}
\renewcommand{\StudyQuestionExtra}{\StudyTarget{MultipleChoice__mt045}\StudyTarget{MultipleChoice__mt045-ridgeomp}}
\StudyQuestion{MT045}{特征选择：零系数、Ridge与OMP}{选择题（多选） / Multiple select}
\input{\StudyRoot/CS5489/exam-review/Midterm/items/MT045/provenance.tex}
\subsection*{English question}
\StudyTarget{MultipleChoice__english-question_23}
\input{\StudyRoot/CS5489/exam-review/Midterm/items/MT045/question-en.tex}
\subsection*{中文题意}
\StudyTarget{MultipleChoice___48}
\input{\StudyRoot/CS5489/exam-review/Midterm/items/MT045/question-zh.tex}
\subsection*{English answer}
\StudyTarget{MultipleChoice__english-answer_23}
\input{\StudyRoot/CS5489/exam-review/Midterm/items/MT045/answer-en.tex}
\subsection*{中文对应答案}
\StudyTarget{MultipleChoice___49}
\input{\StudyRoot/CS5489/exam-review/Midterm/items/MT045/answer-zh.tex}
\input{\StudyRoot/CS5489/exam-review/Midterm/items/MT045/concise-extra.tex}
\renewcommand{\StudyQuestionExtra}{\StudyTarget{MultipleChoice__mt030}\StudyTarget{MultipleChoice__mt030-l1l2}}
\StudyQuestion{MT030}{L1/L2回归的益处与措辞边界}{选择题（多选） / Multiple select}
\input{\StudyRoot/CS5489/exam-review/Midterm/items/MT030/provenance.tex}
\subsection*{English question}
\StudyTarget{MultipleChoice__english-question_24}
\input{\StudyRoot/CS5489/exam-review/Midterm/items/MT030/question-en.tex}
\subsection*{中文题意}
\StudyTarget{MultipleChoice___50}
\input{\StudyRoot/CS5489/exam-review/Midterm/items/MT030/question-zh.tex}
\subsection*{English answer}
\StudyTarget{MultipleChoice__english-answer_24}
\input{\StudyRoot/CS5489/exam-review/Midterm/items/MT030/answer-en.tex}
\subsection*{中文对应答案}
\StudyTarget{MultipleChoice___51}
\input{\StudyRoot/CS5489/exam-review/Midterm/items/MT030/answer-zh.tex}
\input{\StudyRoot/CS5489/exam-review/Midterm/items/MT030/concise-extra.tex}
\renewcommand{\StudyQuestionExtra}{\StudyTarget{MultipleChoice__mt057}\StudyTarget{MultipleChoice__mt057-l1l2elastic-net}}
\StudyQuestion{MT057}{同时使用L1与L2：Elastic Net}{选择题（多选） / Multiple select}
\input{\StudyRoot/CS5489/exam-review/Midterm/items/MT057/provenance.tex}
\subsection*{English question}
\StudyTarget{MultipleChoice__english-question_25}
\input{\StudyRoot/CS5489/exam-review/Midterm/items/MT057/question-en.tex}
\subsection*{中文题意}
\StudyTarget{MultipleChoice___52}
\input{\StudyRoot/CS5489/exam-review/Midterm/items/MT057/question-zh.tex}
\subsection*{English answer}
\StudyTarget{MultipleChoice__english-answer_25}
\input{\StudyRoot/CS5489/exam-review/Midterm/items/MT057/answer-en.tex}
\subsection*{中文对应答案}
\StudyTarget{MultipleChoice___53}
\input{\StudyRoot/CS5489/exam-review/Midterm/items/MT057/answer-zh.tex}
\input{\StudyRoot/CS5489/exam-review/Midterm/items/MT057/concise-extra.tex}
\Needspace{10\baselineskip}\phantomsection\addcontentsline{toc}{chapter}{集成学习}
\renewcommand{\StudyQuestionExtra}{\StudyTarget{MultipleChoice__mt006}\StudyTarget{MultipleChoice__mt006-baggingboosting}}
\StudyQuestion{MT006}{Bagging与Boosting的基本区别}{选择题（多选） / Multiple select}
\input{\StudyRoot/CS5489/exam-review/Midterm/items/MT006/provenance.tex}
\subsection*{English question}
\StudyTarget{MultipleChoice__english-question_26}
\input{\StudyRoot/CS5489/exam-review/Midterm/items/MT006/question-en.tex}
\subsection*{中文题意}
\StudyTarget{MultipleChoice___54}
\input{\StudyRoot/CS5489/exam-review/Midterm/items/MT006/question-zh.tex}
\subsection*{English answer}
\StudyTarget{MultipleChoice__english-answer_26}
\input{\StudyRoot/CS5489/exam-review/Midterm/items/MT006/answer-en.tex}
\subsection*{中文对应答案}
\StudyTarget{MultipleChoice___55}
\input{\StudyRoot/CS5489/exam-review/Midterm/items/MT006/answer-zh.tex}
\input{\StudyRoot/CS5489/exam-review/Midterm/items/MT006/concise-extra.tex}
\renewcommand{\StudyQuestionExtra}{\StudyTarget{MultipleChoice__mt044}\StudyTarget{MultipleChoice__mt044-boosting}}
\StudyQuestion{MT044}{随机森林与Boosting的用途}{选择题（多选） / Multiple select}
\input{\StudyRoot/CS5489/exam-review/Midterm/items/MT044/provenance.tex}
\subsection*{English question}
\StudyTarget{MultipleChoice__english-question_27}
\input{\StudyRoot/CS5489/exam-review/Midterm/items/MT044/question-en.tex}
\subsection*{中文题意}
\StudyTarget{MultipleChoice___56}
\input{\StudyRoot/CS5489/exam-review/Midterm/items/MT044/question-zh.tex}
\subsection*{English answer}
\StudyTarget{MultipleChoice__english-answer_27}
\input{\StudyRoot/CS5489/exam-review/Midterm/items/MT044/answer-en.tex}
\subsection*{中文对应答案}
\StudyTarget{MultipleChoice___57}
\input{\StudyRoot/CS5489/exam-review/Midterm/items/MT044/answer-zh.tex}
\input{\StudyRoot/CS5489/exam-review/Midterm/items/MT044/concise-extra.tex}
\renewcommand{\StudyQuestionExtra}{\StudyTarget{MultipleChoice__mt056}}
\StudyQuestion{MT056}{集成方法：树深、并行与后续轮次}{选择题（多选） / Multiple select}
\input{\StudyRoot/CS5489/exam-review/Midterm/items/MT056/provenance.tex}
\subsection*{English question}
\StudyTarget{MultipleChoice__english-question_28}
\input{\StudyRoot/CS5489/exam-review/Midterm/items/MT056/question-en.tex}
\subsection*{中文题意}
\StudyTarget{MultipleChoice___58}
\input{\StudyRoot/CS5489/exam-review/Midterm/items/MT056/question-zh.tex}
\subsection*{English answer}
\StudyTarget{MultipleChoice__english-answer_28}
\input{\StudyRoot/CS5489/exam-review/Midterm/items/MT056/answer-en.tex}
\subsection*{中文对应答案}
\StudyTarget{MultipleChoice___59}
\input{\StudyRoot/CS5489/exam-review/Midterm/items/MT056/answer-zh.tex}
\input{\StudyRoot/CS5489/exam-review/Midterm/items/MT056/concise-extra.tex}
\Needspace{10\baselineskip}\phantomsection\addcontentsline{toc}{chapter}{历史范围拓展：Gaussian Process}
\renewcommand{\StudyQuestionExtra}{\StudyTarget{MultipleChoice__mt003}\StudyTarget{MultipleChoice__mt003-gaussian-process}}
\StudyQuestion{MT003}{历史拓展：Gaussian Process的假设}{选择题（多选） / Multiple select}
\input{\StudyRoot/CS5489/exam-review/Midterm/items/MT003/provenance.tex}
\subsection*{English question}
\StudyTarget{MultipleChoice__english-question_29}
\input{\StudyRoot/CS5489/exam-review/Midterm/items/MT003/question-en.tex}
\subsection*{中文题意}
\StudyTarget{MultipleChoice___60}
\input{\StudyRoot/CS5489/exam-review/Midterm/items/MT003/question-zh.tex}
\subsection*{English answer}
\StudyTarget{MultipleChoice__english-answer_29}
\input{\StudyRoot/CS5489/exam-review/Midterm/items/MT003/answer-en.tex}
\subsection*{中文对应答案}
\StudyTarget{MultipleChoice___61}
\input{\StudyRoot/CS5489/exam-review/Midterm/items/MT003/answer-zh.tex}
\input{\StudyRoot/CS5489/exam-review/Midterm/items/MT003/concise-extra.tex}
\clearpage\StudyHeading{1}{原卷索引 / Original-paper index}
按原卷题号查找本册位置；不同年份的重复题指向同一MT编号。2021B*保留题纸与配套答案的年份冲突。
\StudyTarget{MultipleChoice__original-paper-index}
\subsection*{2020B期中Quiz}
\StudyTarget{MultipleChoice__2020bquiz}
\begin{longtable}{@{}p{.2\linewidth}p{.56\linewidth}p{.18\linewidth}@{}}
\toprule 原题 & 本册题号与页码 & 分值 \\ \midrule\endhead
Q1 & \StudyRef{MultipleChoice}{mt001}{MT001} & 5 \\
Q2 & \StudyRef{MultipleChoice}{mt002}{MT002} & 5 \\
Q3 & \StudyRef{MultipleChoice}{mt003}{MT003} & 5 \\
Q4 & \StudyRef{MultipleChoice}{mt004}{MT004} & 5 \\
Q5 & \StudyRef{MultipleChoice}{mt005}{MT005} & 5 \\
Q6 & \StudyRef{MultipleChoice}{mt006}{MT006} & 5 \\
\bottomrule\end{longtable}
\subsection*{2021A期中}
\StudyTarget{MultipleChoice__2021a}
\begin{longtable}{@{}p{.2\linewidth}p{.56\linewidth}p{.18\linewidth}@{}}
\toprule 原题 & 本册题号与页码 & 分值 \\ \midrule\endhead
Q1 & \StudyRef{MultipleChoice}{mt014}{MT014} & 5 \\
Q2 & \StudyRef{MultipleChoice}{mt015}{MT015} & 5 \\
Q3 & \StudyRef{MultipleChoice}{mt016}{MT016} & 5 \\
Q4 & \StudyRef{MultipleChoice}{mt017}{MT017} & 5 \\
Q5 & \StudyRef{MultipleChoice}{mt018}{MT018} & 5 \\
Q6 & \StudyRef{MultipleChoice}{mt019}{MT019} & 5 \\
\bottomrule\end{longtable}
\subsection*{2021B*期中}
\StudyTarget{MultipleChoice__2021b}
\begin{longtable}{@{}p{.2\linewidth}p{.56\linewidth}p{.18\linewidth}@{}}
\toprule 原题 & 本册题号与页码 & 分值 \\ \midrule\endhead
Q1 & \StudyRef{MultipleChoice}{mt027}{MT027} & 5 \\
Q2 & \StudyRef{MultipleChoice}{mt028}{MT028} & 5 \\
Q3 & \StudyRef{MultipleChoice}{mt029}{MT029} & 5 \\
Q4 & \StudyRef{MultipleChoice}{mt030}{MT030} & 5 \\
Q5 & \StudyRef{MultipleChoice}{mt031}{MT031} & 5 \\
Q6 & \StudyRef{MultipleChoice}{mt032}{MT032} & 5 \\
\bottomrule\end{longtable}
\subsection*{2023A期中}
\StudyTarget{MultipleChoice__2023a}
\begin{longtable}{@{}p{.2\linewidth}p{.56\linewidth}p{.18\linewidth}@{}}
\toprule 原题 & 本册题号与页码 & 分值 \\ \midrule\endhead
Q1 & \StudyRef{MultipleChoice}{mt040}{MT040} & 5 \\
Q2 & \StudyRef{MultipleChoice}{mt041}{MT041} & 5 \\
Q3 & \StudyRef{MultipleChoice}{mt042}{MT042} & 5 \\
Q4 & \StudyRef{MultipleChoice}{mt043}{MT043} & 5 \\
Q5 & \StudyRef{MultipleChoice}{mt044}{MT044} & 5 \\
Q6 & \StudyRef{MultipleChoice}{mt045}{MT045} & 5 \\
\bottomrule\end{longtable}
\subsection*{2023B期中}
\StudyTarget{MultipleChoice__2023b}
\begin{longtable}{@{}p{.2\linewidth}p{.56\linewidth}p{.18\linewidth}@{}}
\toprule 原题 & 本册题号与页码 & 分值 \\ \midrule\endhead
Q1 & \StudyRef{MultipleChoice}{mt052}{MT052} & 5 \\
Q2 & \StudyRef{MultipleChoice}{mt053}{MT053} & 5 \\
Q3 & \StudyRef{MultipleChoice}{mt054}{MT054} & 5 \\
Q4 & \StudyRef{MultipleChoice}{mt055}{MT055} & 5 \\
Q5 & \StudyRef{MultipleChoice}{mt056}{MT056} & 5 \\
Q6 & \StudyRef{MultipleChoice}{mt057}{MT057} & 5 \\
\bottomrule\end{longtable}
\subsection*{2025A期中}
\StudyTarget{MultipleChoice__2025a}
\begin{longtable}{@{}p{.2\linewidth}p{.56\linewidth}p{.18\linewidth}@{}}
\toprule 原题 & 本册题号与页码 & 分值 \\ \midrule\endhead
Q1 & \StudyRef{MultipleChoice}{mt040}{MT040} & 5 \\
Q2 & \StudyRef{MultipleChoice}{mt041}{MT041} & 5 \\
Q3 & \StudyRef{MultipleChoice}{mt042}{MT042} & 5 \\
Q4 & \StudyRef{MultipleChoice}{mt043}{MT043} & 5 \\
Q5 & \StudyRef{MultipleChoice}{mt044}{MT044} & 5 \\
Q6 & \StudyRef{MultipleChoice}{mt057}{MT057} & 5 \\
\bottomrule\end{longtable}
\end{document}
